% from ega4-18-fr.tex
\oldpage[IV-4]{182}

En vertu de la relation $k(y_1)=k(y)$, la fibre $f_1^{-1}(y_1)$ dans $S_1$ est isomorphe à $f^{-1}(y)$, donc $X_1\cap f_1^{-1}(y_1)$ est fini, discret et a $n-1$ points. Comme $f_1|X_1$ est localement de type fini et séparé, on applique l'hypothèse de récurrence à ce morphisme: il y a un morphisme étale $Y'\to Y_1$, un point $y'\in Y'$ au-dessus de $y_1$ tel que $k(y')=k(y_1)$, et, si $X'=X\times_Y Y'$ et si $p:X'\to S_1$ est la projection canonique, on a une décomposition de $p^{-1}(X_1)$ en somme de deux sous-préschémas $U'$ et $X'_2$ induits sur des ouverts de $X'$ tels que $X'_2\cap f'^{-1}(y')=\varnothing$ et que $U'$ soit fini sur $Y'$. Par ailleurs $V'=p^{-1}(V_1)$ est fini sur $Y'$ et $X'$ est somme de $U'$, $V'$ et $X'_2$; on répondra donc à la question en prenant $X'_1$ somme de $U'$ et $V'$.

Le corollaire suivant améliore (8.11.1):

\begin{corollary}[18.12.4]
\label{IV.18.12.4}
Pour qu'un morphisme $f:X\to Y$ soit fini, il faut et il suffit qu'il soit séparé, universellement fermé, localement de type fini et que, pour tout $y\in Y$, $f^{-1}(y)$ soit un espace fini discret; en particulier, un morphisme propre et quasi-fini est fini.
\end{corollary}

En effet, on peut appliquer le corollaire (18.12.3) à un point quelconque $y$ de $Y$. Avec les notations de ce corollaire, $f'$ est un morphisme fermé, donc, puisque $X'_2$ est fermé dans $X'$, $f'(X'_2)$ est fermé dans $Y'$ et ne contient pas $y'$; il existe par suite un voisinage ouvert $U'$ de $y'$ dans $Y'$ tel que $U'\to Y$ soit de type fini et que $f'^{-1}(U')=X\times_Y U'$ soit fini sur $U'$. Soit $U$ l'image de $U'$ dans $Y$, qui est un ouvert de $Y$ en vertu de (II, 3.1) puisque le morphisme $Y'\to Y$ est étale, donc plat et localement de présentation finie; on a évidemment encore $X\times_Y U'=X\times_U U'=f^{-1}(U)\times_U U'$. Puisque maintenant le morphisme $U'\to U$ est fidèlement plat et quasi-compact, on déduit de (2.7.1, (xv)) que le morphisme $f^{-1}(U)\to U$, restriction de $f$, est fini, ce qui prouve le corollaire.

\begin{remark}[18.12.5]
\label{IV.18.12.5}
On notera que dans la démonstration de (18.12.4) on n'a pas utilisé l'hypothèse que $f$ est universellement fermé sous sa forme générale, mais seulement que $f_{(Y')}$ est un morphisme fermé pour tout morphisme étale $Y'\to Y$.
\end{remark}

\begin{corollary}[18.12.6]
\label{IV.18.12.6}
Soit $f:X\to Y$ un morphisme de préschémas. Les conditions suivantes sont équivalentes:
\begin{enumerate}
  \item[$a)$] $f$ est une immersion fermée.
  \item[$b)$] $f$ est un monomorphisme propre.
  \item[$c)$] $f$ est propre et, pour tout $y\in Y$, le $k(y)$-préschéma $X_y=f^{-1}(y)$ est radiciel et géométriquement réduit sur $k(y)$ (c'est-à-dire vide ou $k(y)$-isomorphe à $\Spec(k(y))$).
\end{enumerate}
\end{corollary}

Cela se déduit de (18.12.4) comme (8.11.5) se déduit de (8.11.1).

\begin{proposition}[18.12.7]
\label{IV.18.12.7}
Soient $f:X\to Y$ un morphisme localement de type fini, $y$ un point de $Y$. Pour qu'il existe un voisinage ouvert $U$ de $y$ tel que la restriction $f^{-1}(U)\to U$ de $f$ soit une immersion fermée, il faut et il suffit que $X_y$ soit un $k(y)$-préschéma radiciel et géométriquement réduit sur $k(y)$ (c'est-à-dire vide ou $k(y)$-isomorphe à $\Spec(k(y))$) et qu'il existe un voisinage ouvert $V$ de $y$ tel que la restriction $f^{-1}(V)\to V$ de $f$ soit un morphisme universellement fermé.
\end{proposition}

Les conditions étant évidemment nécessaires, prouvons qu'elles sont suffisantes. On peut se borner au cas où $X_y\neq\varnothing$. En vertu de la seconde condition, on peut déjà supposer $f$ universellement fermé. Montrons en outre qu'on peut supposer $f$ affine, donc séparé; cela résulte du lemme suivant:

\begin{lemma}[18.12.7.1]
\label{IV.18.12.7.1}
Soient $f:X\to Y$ un morphisme fermé, $y\in Y$ un point tel que tout voisinage de $X_y$ contienne un voisinage affine de $X_y$ (condition toujours vérifiée lorsque $X_y$ est vide ou réduite à un seul point). Alors il existe un voisinage ouvert affine $U$ de $y$ tel que la restriction $f^{-1}(U)\to U$ de $f$ soit un morphisme affine.
\end{lemma}

En effet, soit $U_0$ un voisinage ouvert affine de $y$ dans $Y$ et soit $V$ un voisinage ouvert affine de $X_y$ contenu dans $f^{-1}(U_0)$. Puisque $f$ est fermé, il existe un voisinage ouvert affine $U\subset U_0$ de $y$ tel que $f^{-1}(U)\subset V$. Comme la restriction $g:V\to U_0$ de $f$ est un morphisme affine, il en est de même de sa restriction $f^{-1}(U)\to U$ (II, 1.2.5).

Supposons donc $f$ affine et universellement fermé, et montrons qu'il existe un voisinage ouvert $U$ de $y$ tel que la restriction $f^{-1}(U)\to U$ soit un morphisme fini; comme on peut supposer $Y$ et $X$ affines, on peut supposer $f$ de type fini, donc propre, et il suffit de prouver qu'il existe un voisinage $U$ de $y$ tel que $f^{-1}(U)\to U$ soit un morphisme quasi-fini (18.12.4). Mais il existe un voisinage $V$ de l'ensemble $X_y$ (réduit à un seul point) tel que $f|V$ soit quasi-fini (13.1.4), et comme $f$ est fermé il y a un voisinage $U$ de $y$ tel que $f^{-1}(U)\subset V$, ce qui achève la démonstration, compte tenu de (18.12.6).

\begin{proposition}[18.12.8]
\label{IV.18.12.8}
Soit $f:X\to Y$ un morphisme de préschémas. Pour que $f$ soit entier, il faut et il suffit que $f$ soit affine et universellement fermé.
\end{proposition}

Les conditions sont nécessaires (II, 6.1.10); prouvons qu'elles sont suffisantes. On peut supposer $Y=\Spec(A)$ affine, donc aussi $X=\Spec(B)$, et il faut prouver que tout élément $b\in B$ est entier sur $A$ (II, 6.1.1). Soit $B'$ la sous-$A$-algèbre de $B$ engendrée par $b$, qui est une $A$-algèbre de type fini; posons $X'=\Spec(B')$, de sorte que $f:X\to Y$ se factorise en
\[
  X\xrightarrow{g}X'\xrightarrow{h}Y,
\]
où $h$ est de type fini, et $g$ est dominant puisque l'homomorphisme $B'\to B$ est injectif (I, 1.2.7). Comme $h$ est séparé et $h\circ g$ universellement fermé, $g$ est universellement fermé (II, 5.4.3 et

